\makeabstract
    {
    Ghost Hunting: Tracking Celestial Objects
    }
    {
    Hermione Arias Soto\textsuperscript{1}, Meadow Perez\textsuperscript{1}
    }
    {
    \textsuperscript{1} Department of Physics \& Astronomy, Amherst College, Amherst, MA
    }
    {
    Astronomers know a lot about how stars form and hypothesize that the formation of substellar objects like brown dwarfs and planets could be similar. We aim to better understand the formation of substellar objects by analyzing how quickly they grow (their mass accretion rates) in comparison to their mass. 
    }
    {
    To observe the relationship between mass and accretion rate, we expanded the Comprehensive Archive of Substellar and Planetary Accretion Rates (CASPAR). We ingest over 300 new entries from published literature by documenting accretion tracers, equivalent widths, mass, radius and other variables necessary to determine accretion rates. To maintain consistency across all objects, we uniformly rederived these values using standardized scaling relations.
    }
    {
     Our new accretion values are similar to the slope established by previous studies which further suggests that brown dwarfs and planets form differently than stars.
    }
    {
    The integration of over 300 new literature values highlights that there are significant gaps in our understanding of the formation of substellar objects. By continuing to expand the database, we will be able to constrain theoretical models of substellar evolution and better understand how our own solar system formed.
    }

    \newpage
    